
\section*{Abstract}
\addcontentsline{toc}{section}{Abstract}
\vspace{2cm}

\large
This Progress Report summarises progress on \textit{GiTrip}, a web-based trip planning system that applies \textit{Git}-style version control concepts (repositories, commits, branches, merges, and conflict resolution) to itinerary planning. The project targets a common problem in group travel: plans evolve quickly, collaborators make parallel edits, and typical tools either optimise a single route (e.g., map applications) or support collaboration without structured change management (e.g., generic documents). \textit{GiTrip} combines version control semantics with a constraint-based scheduler that considers travel time, daily active hours, per-stop stay durations, opening-hours constraints, and user preferences (compact vs.\ sparse timing, time-of-day focus, strict ordering, and transport mode).

By the end of Term~1, the project has delivered a working vertical slice: repositories and history in SQLite; a planner that produces feasible day-by-day schedules; server-side routing integration using Openrouteservice (matrix and per-leg directions) and Google Directions for transit; place search via a Nominatim proxy; interactive visualisation through an SVG timeline editor and Leaflet map; and a merge UI with conflict detection and resolution. Robust fallbacks are implemented so the system remains functional under API quota limits or outages.

The remainder of the project focuses on deployment and Progressive Web App support, security hardening, systematic testing and evaluation (including a UAT dry-run), and final report/viva preparation. An updated timetable and risk plan are provided to demonstrate feasibility for successful completion.

\vspace{1cm}

\noindent \textit{Keywords: Flux Capacitor, 1.21 Gigawatts, Calvin Klein}